\documentclass[../../main.tex]{subfiles}
\graphicspath{{../../image/}{../../figure/}}
\begin{document}

\section{阻尼的重分布}\label{section:ch16-阻尼重分布}

由于矩阵 $A$ 的对角占优\eqref{align:ch16-s02-e04}对预期应用而言过于苛刻，我们在此寻求仍能为 Cauchy 问题\eqref{align:ch16-s02-e01}、\eqref{align:ch16-s01-e02}在大范围诱导 $BV$ 解存在性的替代假设。

对 $\omega>0$，记 $\mathscr{C}_\omega$ 为初始数据 $U_0$ 的类，具有如下性质：\eqref{align:ch16-s02-e01}、\eqref{align:ch16-s01-e02}的任一可容 $BV$ 解 $U$ 在其存在的任意时间区间上都满足
\begin{align}
\int_{-\infty}^{\infty}|U(x,t)|\,\mathrm{d}x\leqslant\omega,
\label{align:ch16-s03-e01}
\end{align}
目标是证明：当初值取在 $\mathscr{C}_\omega$ 中时，\eqref{align:ch16-s02-e01}的可容 $BV$ 解在大范围存在，只要矩阵 $A$ 的主对角线上的元素为正：
\begin{align}
A_{ii}>0,\qquad i=1,\ldots,n.
\label{align:ch16-s03-e02}
\end{align}
正如我们将在\cref{section:ch16-原点锥奇异熵l1稳定}与\cref{section:ch16-部分耗散源l1稳定}中看到的，\eqref{align:ch16-s03-e02}是刻画松弛现象的系统的典型性质。此外，在此类系统中，类 $\mathscr{C}_\omega$ 足够丰富，能包含平衡态附近令人感兴趣的初始数据。

\begin{theorem}\label{theorem:ch16-1631}
考虑齐次、严格双曲的平衡律组\eqref{align:ch16-s02-e01}，其特征族要么真非线性、要么线性退化。设 $0$ 为平衡态，$P(0)=0$，并假设由\eqref{align:ch16-s02-e03}定义的矩阵 $A$ 的元素 $A_{ii}$ 为正\eqref{align:ch16-s03-e02}。则存在正数 $\omega_0$ 与 $\delta_0$，使得对任意属于 $\mathscr{C}_\omega$ 的初始数据 $U_0$（其中 $\omega<\omega_0$）并满足\eqref{align:ch16-s01-e10}（其中 $\delta<\delta_0$），Cauchy 问题\eqref{align:ch16-s02-e01}、\eqref{align:ch16-s01-e02}在 $(-\infty,\infty)\times[0,\infty)$ 上拥有一个可容 $BV$ 解 $U$，具有性质
\begin{align}
TV_{(-\infty,\infty)}U(\cdot,t)\leqslant c_0\omega+c_1\delta e^{-\nu t},\qquad 0\leqslant t<\infty,
\label{align:ch16-s03-e03}
\end{align}
其中 $c_0$、$c_1$ 与 $\nu$ 为正常数，与 $U_0$ 无关。此外，若
\begin{align}
\int_{-\infty}^{\infty}|U(x,t)|\,\mathrm{d}x\to0,\qquad\text{当 }t\to\infty,
\label{align:ch16-s03-e04}
\end{align}
则
\begin{align}
TV_{(-\infty,\infty)}U(\cdot,t)\to0,\qquad\text{当 }t\to\infty.
\label{align:ch16-s03-e05}
\end{align}
\end{theorem}

我们在此给出上述定理证明的路线图，技术细节推迟到\cref{section:ch16-变差界}。

在控制松弛现象的系统中，矩阵 $A$ 上的对角占优条件\eqref{align:ch16-s02-e04}并不成立，尽管源是耗散的——因为阻尼作用在系统的各方程间分布不均。这一缺陷将在此通过状态向量的变换、在更公平的基础上重分布阻尼来弥补，如下所述。

假设 $U$ 是\eqref{align:ch16-s02-e01}、\eqref{align:ch16-s01-e02}的一个可容 $BV$ 解，定义在带形 $(-\infty,\infty)\times[0,T)$ 上，且对所有 $t\in[0,T)$，$U(\cdot,t)$ 在 $(-\infty,\infty)$ 上可积。我们引入函数
\begin{align}
\Phi(x,t)=\int_{-\infty}^x NU(y,t)\,\mathrm{d}y,
\label{align:ch16-s03-e06}\\
Z(x,t)=\int_{-\infty}^x NP(U(y,t))\,\mathrm{d}y,
\label{align:ch16-s03-e07}
\end{align}
其中 $N$ 是一个待在下面指定的 $n\times n$ 矩阵。我们注意到 $\Phi$ 是 Lipschitz 的，且
\begin{align}
\partial_x\Phi(x,t)=NU(x,t),\qquad\partial_t\Phi(x,t)=-NF(U(x,t))-Z(x,t).
\label{align:ch16-s03-e08}
\end{align}
现在我们以新的状态向量
\begin{align}
\hat U=U-\Phi
\label{align:ch16-s03-e09}
\end{align}
代替 $U$，并把\eqref{align:ch16-s02-e01}改写成关于 $\hat U$ 的系统，即
\begin{align}
\partial_t\hat U(x,t)+\partial_x\hat F(\hat U(x,t),\Phi(x,t))+\hat P(\hat U(x,t),\Phi(x,t),Z(x,t))=0,
\label{align:ch16-s03-e10}
\end{align}
其中
\begin{align}
\hat F(\hat U,\Phi)=F(\hat U+\Phi)-F(\Phi),
\label{align:ch16-s03-e11}\\
\hat P(\hat U,\Phi,Z)=P(\hat U+\Phi)-NF(\hat U+\Phi)+DF(\Phi)N[\hat U+\Phi]-Z.
\label{align:ch16-s03-e12}
\end{align}

从相对简单的\eqref{align:ch16-s02-e01}切换到笨重的\eqref{align:ch16-s03-e10}（它甚至不是一个封闭系统）的动机是：若我们假定 $\Phi$ 与 $Z$ 以某种方式已知，并把\eqref{align:ch16-s03-e10}视为关于 $\hat U$ 的非齐次平衡律组，则在\eqref{align:ch16-s02-e03}的位置上，我们得到矩阵 $\hat A$，其元素为
\begin{align}
\hat A_{ij}=L_j(0)D\hat P(0,0,0)R_j(0)=A_{ij}+[\lambda_i(0)-\lambda_j(0)]\Delta_{ij},
\label{align:ch16-s03-e13}
\end{align}
其中 $\Delta=L(0)NR(0)$。这提供了一个机会：通过适当选取矩阵 $N$，使 $\hat A$ 对角占优。特别地，取 $N=R(0)\Delta L(0)$，其中 $\Delta_{ii}=0$（$i=1,\ldots,n$），且
\begin{align}
\Delta_{ij}=\frac{A_{ij}}{\lambda_i(0)-\lambda_j(0)},\qquad i\neq j,
\label{align:ch16-s03-e14}
\end{align}
则 $\hat A$ 为对角矩阵，$\hat A_{ii}=A_{ii}$（$i=1,\ldots,n$），$\hat A_{ij}=0$（$i\neq j$）。在这种情形下，当\eqref{align:ch16-s03-e02}成立时，$\hat A$ 对角占优。

上述考虑提示了证明\cref{theorem:ch16-1631}的如下程序。从位于 $\mathscr{C}_\omega$ 且满足\eqref{align:ch16-s01-e10}的 $U_0$ 出发，当 $\omega$ 与 $\delta$ 充分小时，\cref{theorem:ch16-1611}保证在某个时间区间 $[0,T)$ 上存在 Cauchy 问题\eqref{align:ch16-s02-e01}、\eqref{align:ch16-s01-e02}的局部可容 $BV$ 解 $U$。事实上，已知（见\cref{section:ch16-注记}所引文献）：类似于第 14.9 节中讨论的守恒律组情形，形式如\eqref{align:ch16-s02-e01}的齐次平衡律组生成标准 Riemann 半群，从而诱导解 $U$ 的唯一性。

从解 $U$ 出发，我们可以通过\eqref{align:ch16-s03-e06}与\eqref{align:ch16-s03-e07}在 $(-\infty,\infty)\times[0,T)$ 上确定 $\Phi$ 与 $Z$，然后在初始条件 $\hat U(x,0)=\hat U_0(x)$ 下求解\eqref{align:ch16-s03-e10}的 Cauchy 问题来构造 $\hat U$，其中
\begin{align}
\hat U_0(x)=U_0(x)-\int_{-\infty}^x NU_0(y)\,\mathrm{d}y,\qquad-\infty<x<\infty.
\label{align:ch16-s03-e15}
\end{align}
最后，我们可以把 $U$ 恢复为 $\hat U+\Phi$。收益在于：源 $\hat P$ 上的耗散作用（体现为 $\hat A$ 的对角占优）将通过从 $\hat U$ 重构 $U$ 而显现，诱导出关于变差的估计\eqref{align:ch16-s03-e03}，其中常数 $c_0$、$c_1$ 与 $\mu$ 与 $T$ 无关。经由上述程序推导\eqref{align:ch16-s03-e03}的细节将在下一节\cref{section:ch16-变差界}中给出。

有了估计\eqref{align:ch16-s03-e03}，我们可以立即把解 $U$ 延拓到整个上半平面 $(-\infty,\infty)\times[0,\infty)$。此外，在\eqref{align:ch16-s03-e03}存在的情况下，\eqref{align:ch16-s03-e04}蕴涵\eqref{align:ch16-s03-e05}。为看到这一点，假设\eqref{align:ch16-s03-e04}成立，固定任意 $\varepsilon>0$，令 $\bar\omega=\varepsilon/2c_0$，并确定 $\bar t>0$ 使得当 $t\in[\bar t,\infty)$ 时 $\|U(\cdot,t)\|_{L^1}<\bar\omega$。于是 $U(\cdot,\bar t)\in\mathscr{C}_{\bar\omega}$。此外，由\eqref{align:ch16-s03-e03}，$TV_{(-\infty,\infty)}U(\cdot,\bar t)<c_0\omega+c_1\delta$。接着我们把时间原点从 $0$ 平移到 $\bar t$ 后应用\eqref{align:ch16-s03-e03}，得到
\begin{align}
TV_{(-\infty,\infty)}U(\cdot,t)\leqslant c_0\bar\omega+c_1(c_0\omega+c_1\delta)e^{-\nu(t-\bar t)},\qquad\bar t\leqslant t<\infty.
\label{align:ch16-s03-e16}
\end{align}
现在显然存在 $t_\varepsilon>\bar t$，使得对所有 $t\geqslant t_\varepsilon$，$TV_{(-\infty,\infty)}U(\cdot,t)<\varepsilon$。

\end{document}
